% NEJM Case 4-2026: 80-Year-Old Woman with Cough and Hypoxemia
% HIV-AIDS with Pneumocystis jirovecii and Cryptococcus pneumonia
% 作者: 謝慕揚, MD, PhD, FESC

\documentclass[12pt,a4paper]{article}

% Drake_preamble.tex - LaTeX Preamble Template
% 作者: 謝慕揚, MD, PhD, FESC
% 
% 使用說明:
% 1. 表格請使用 longtable，不要有垂直分隔線 |
% 2. 表格內換行使用 \newline，不要使用 \\
% 3. 藥物名稱、檢驗、檢查請維持英文
% 4. Reference 格式請使用 \href 加上驗證過的 DOI

% Including essential packages for Chinese typesetting
\usepackage{xeCJK}
\usepackage{fontspec}
% Configuring Chinese font with Noto Serif CJK TC, fallback to SimSun
\IfFontExistsTF{Noto Serif CJK TC}{
  \setCJKmainfont{Noto Serif CJK TC}
  \setCJKsansfont{Noto Serif CJK TC}
  \setCJKmonofont{Noto Serif CJK TC}
}{\setCJKmainfont{SimSun}
  \setCJKsansfont{SimSun}
  \setCJKmonofont{SimSun}
}

% Including other necessary packages
\usepackage{geometry}
\usepackage{titlesec}
\usepackage{enumitem}
\usepackage{xcolor}
\usepackage{colortbl}
\usepackage{tcolorbox}
\usepackage{booktabs}
\usepackage{longtable}
\usepackage{array}
\usepackage{graphicx}
\usepackage{tikz}
\usepackage{fancyhdr}
\usepackage{multicol}
\usepackage{datetime2}
\usepackage{hyperref}
\usepackage{mdframed}
\usepackage{amsmath}
\usepackage{amssymb}
\usepackage{multirow}

\hypersetup{
    colorlinks=true,
    linkcolor=emergency_blue,
    citecolor=emergency_blue,
    urlcolor=emergency_blue,
    pdfborder={0 0 0}
}

% Defining page geometry
\geometry{
  top=2.5cm,
  bottom=2.5cm,
  left=2cm,
  right=2cm,
  headheight=25pt
}

% Adjusting topmargin to compensate for increased headheight
\addtolength{\topmargin}{-10pt}

% Defining colors
\definecolor{emergency_red}{RGB}{186,24,27}
\definecolor{emergency_blue}{RGB}{0,114,188}
\definecolor{emergency_gray}{RGB}{120,120,120}
\definecolor{highlight_yellow}{RGB}{255,250,205}

% Setting title formats
\titleformat{\section}[block]{\Large\bfseries\color{emergency_red}}{\thesection}{10pt}{}
\titleformat{\subsection}[block]{\large\bfseries\color{emergency_blue}}{\thesubsection}{8pt}{}
\titleformat{\subsubsection}[block]{\normalsize\bfseries\color{emergency_gray}}{\thesubsubsection}{6pt}{}

% Defining custom environment for important notes
\newmdenv[
    linecolor=emergency_red,
    backgroundcolor=highlight_yellow,
    linewidth=2pt,
    topline=true,
    bottomline=true,
    leftline=true,
    rightline=true,
    innertopmargin=10pt,
    innerbottommargin=10pt,
    innerleftmargin=10pt,
    innerrightmargin=10pt
]{importantbox}

% Defining note command with tcolorbox
\newcommand{\note}[1]{
    \par
    \noindent
    \begin{tcolorbox}[
        colback=emergency_blue!5!white,
        colframe=emergency_blue,
        title=\textbf{重點提醒},
        fonttitle=\bfseries,
        boxrule=0.5pt,
        arc=3pt,
        left=6pt,
        right=6pt,
        top=6pt,
        bottom=6pt
    ]
    #1
    \end{tcolorbox}
}

% Key findings box
\newcommand{\keyfinding}[1]{
    \par
    \noindent
    \begin{tcolorbox}[
        colback=yellow!10!white,
        colframe=orange!75!black,
        title=\textbf{核心發現摘要},
        fonttitle=\bfseries,
        boxrule=0.5pt,
        arc=3pt
    ]
    #1
    \end{tcolorbox}
}

% Defining custom date format for ISO 8601
\DTMnewdatestyle{iso}{
  \renewcommand{\DTMdisplaydate}[4]{\number##1-\number##2-\number##3}
  \renewcommand{\DTMDisplaydate}[4]{\number##1-\number##2-\number##3}
}
\DTMsetdatestyle{iso}
\newcommand{\mydate}{\DTMtoday}


% Custom header for this document
\pagestyle{fancy}
\fancyhf{}
\fancyhead[L]{\textcolor{emergency_red}{NEJM Case Records 教學講義}}
\fancyhead[R]{\textcolor{emergency_blue}{HIV/AIDS \& PJP}}
\fancyfoot[C]{\textcolor{emergency_gray}{\thepage}}
\renewcommand{\headrulewidth}{1pt}
\renewcommand{\headrule}{\hbox to\headwidth{\color{emergency_red}\leaders\hrule height \headrulewidth\hfill}}

\begin{document}

%============================================================
% Title Page
%============================================================
\begin{center}
{\LARGE\bfseries\textcolor{emergency_red}{NEJM Case Records of the Massachusetts General Hospital}}\\[0.5cm]
{\Large\textbf{Case 4-2026: An 80-Year-Old Woman with Cough and Hypoxemia}}\\[0.3cm]
{\large 80歲女性，咳嗽合併低血氧}\\[0.5cm]
{\normalsize \href{https://doi.org/10.1056/NEJMcpc2513543}{N Engl J Med 2026;394:498-507. DOI: 10.1056/NEJMcpc2513543}}\\[0.3cm]
{\normalsize 整理：謝慕揚 MD, PhD, FESC \quad 日期：\mydate}
\end{center}

\vspace{0.5cm}

%============================================================
\section{病例摘要}
%============================================================

\subsection{主訴與現病史}
\begin{itemize}
    \item 80歲女性，因咳嗽及低血氧入院
    \item 發病前身體狀況良好，可騎腳踏車外出活動
    \item \textbf{8週前}：出現productive cough、rhinorrhea、headache、fatigue、poor appetite
    \item \textbf{5週前}：左側腰部出現帶狀皰疹 (herpes zoster)，接受valacyclovir治療14天
    \item \textbf{6天前}：出現主觀發燒，至診所就醫
    \begin{itemize}
        \item 體重57.4 kg，9個月內減輕6.1 kg
        \item SpO$_2$ 94\% (room air)
        \item 胸部X光：肺部正常，fissure處可能有少量積液
        \item 處方doxycycline 5天療程
    \end{itemize}
    \item \textbf{本次入院}：完成doxycycline療程後隔天出現呼吸困難
\end{itemize}

\subsection{過去病史}
\begin{itemize}
    \item Allergic rhinitis、anxiety、chronic kidney disease
    \item \textbf{1年前}：診斷metastatic squamous-cell carcinoma of unknown primary site
    \begin{itemize}
        \item 左腋下淋巴結切除，另7顆淋巴結無癌細胞
        \item PET-CT：左頸部淋巴結及左扁桃體區代謝增加
        \item 左扁桃體切除後未發現癌症
    \end{itemize}
\end{itemize}

\subsection{旅遊史與社會史}
\begin{itemize}
    \item 曾至南美洲、中美洲、東非旅行
    \item \textbf{6年前}最近一次國際旅行：肯亞、盧旺達、坦尚尼亞
    \item 2年前仍有性行為
    \item 終生不吸菸，不使用毒品，每週飲酒一杯
\end{itemize}

\subsection{入院身體檢查}
\begin{itemize}
    \item Temperature 37.3°C, BP 158/86 mmHg, HR 93 bpm, RR 16/min
    \item \textbf{SpO$_2$ 83\% (room air) → 93\% with O$_2$ 4 L/min via nasal cannula}
    \item 肺部聽診：雙側rhonchi，右下肺葉最明顯
    \item 心臟檢查：規則心律，無雜音
    \item 無clubbing或水腫
\end{itemize}

%============================================================
\section{檢驗數據}
%============================================================

\begin{longtable}{p{6cm}p{4cm}p{4cm}}
\toprule
\textbf{檢驗項目} & \textbf{結果} & \textbf{參考值} \\
\midrule
\endfirsthead
\toprule
\textbf{檢驗項目} & \textbf{結果} & \textbf{參考值} \\
\midrule
\endhead
\bottomrule
\endfoot
Hematocrit (\%) & 35.3 & 36--46 \\
Hemoglobin (g/dL) & 12.1 & 12--16 \\
WBC (per $\mu$L) & 7810 & 4000--11000 \\
\quad Neutrophils & 6230 & 1920--7600 \\
\quad Lymphocytes & \textbf{740} $\downarrow$ & 720--4100 \\
\quad Monocytes & 690 & 160--1100 \\
Platelet (per $\mu$L) & 296,000 & 150,000--450,000 \\
\midrule
\textbf{C-reactive protein (mg/L)} & \textbf{52.7} $\uparrow$ & <8.0 \\
\textbf{ESR (mm/hr)} & \textbf{58} $\uparrow$ & 0--29 \\
Procalcitonin (ng/mL) & 0.11 & <0.25 \\
\midrule
\multicolumn{3}{l}{\textit{Venous blood gas}} \\
pH & 7.48 & 7.30--7.40 \\
PO$_2$ (mmHg) & 49 & 35--50 \\
PCO$_2$ (mmHg) & 32 & 38--50 \\
\end{longtable}

\note{Procalcitonin正常提示非典型細菌感染可能性較低，需考慮其他病因。Lymphopenia合併發炎指標上升是重要線索。}

%============================================================
\section{影像學發現}
%============================================================

\subsection{胸部X光}
\begin{itemize}
    \item Hazy perihilar opacities，血管紋路模糊
    \item 肺底部線狀陰影，可能為subsegmental atelectasis
    \item 心臟輪廓正常，無pleural effusion
\end{itemize}

\subsection{胸部CT特徵}
\begin{itemize}
    \item \textbf{Diffuse ground-glass opacities (GGO)}：雙側肺部
    \item \textbf{Central distribution}（非周邊分布）
    \item \textbf{Apical predominance}（肺尖為主，肺底相對保留）
    \item Trace interlobular septal thickening
    \item 無pleural effusion、pulmonary nodules、cysts、bronchiectasis
    \item 無lymphadenopathy
\end{itemize}

\keyfinding{
CT影像特徵 -- \textbf{Central ground-glass opacities with apical predominance} -- 是\textit{Pneumocystis jirovecii} pneumonia (PJP) 的典型表現！
}

%============================================================
\section{鑑別診斷思路}
%============================================================

\subsection{感染性疾病}

\subsubsection{細菌性肺炎}
\begin{itemize}
    \item 症狀進展8週，不符合典型細菌性肺炎病程
    \item Atypical pneumonia已排除（\textit{C. pneumoniae}、\textit{M. pneumoniae}陰性）
    \item Doxycycline及azithromycin治療無效
    \item Procalcitonin正常
\end{itemize}

\subsubsection{病毒性肺炎}
\begin{itemize}
    \item SARS-CoV-2、influenza、RSV及extended viral panel均陰性
    \item Disseminated VZV pneumonia？
    \begin{itemize}
        \item 通常在皮疹出現後1週內發生，本案為5週後
        \item VZV pneumonitis典型表現為diffuse nodular opacities
    \end{itemize}
\end{itemize}

\subsection{腫瘤相關}
\begin{itemize}
    \item 病人有squamous cell carcinoma病史
    \item 但已完成lymphadenectomy，無殘餘疾病
    \item CT無lymphadenopathy、mass、nodular septal thickening
    \item 不符合lymphoma或lymphangitic carcinomatosis表現
\end{itemize}

\subsection{自體免疫與結締組織疾病}
\begin{itemize}
    \item ANCA-associated vasculitis可造成GGO（肺出血）
    \item 但通常有多器官侵犯（腎絲球腎炎、神經病變、紫斑）
    \item 本案無multisystem involvement，不支持此診斷
\end{itemize}

\subsection{免疫缺損宿主的考量}

\note{
\textbf{關鍵線索}：
\begin{enumerate}
    \item 症狀緩慢進展8週
    \item 最近診斷herpes zoster（可能暗示免疫缺損）
    \item 曾至HIV高盛行率地區旅行（東非）
    \item 2年前仍有性行為
\end{enumerate}
}

\subsubsection{HIV感染相關伺機性感染}
依CD4 cell count可能出現的伺機性感染：
\begin{longtable}{p{5cm}p{9cm}}
\toprule
\textbf{CD4 Count} & \textbf{可能的伺機性感染} \\
\midrule
\endfirsthead
\toprule
\textbf{CD4 Count} & \textbf{可能的伺機性感染} \\
\midrule
\endhead
\bottomrule
\endfoot
Any count & Tuberculosis, VZV reactivation \\
$\leq$250/$\mu$L & Coccidioidomycosis \\
$\leq$200/$\mu$L & \textbf{\textit{Pneumocystis jirovecii} pneumonia (PJP)} \\
$\leq$150/$\mu$L & Histoplasmosis \\
$\leq$100/$\mu$L & Cryptococcus, CMV \\
$\leq$50/$\mu$L & \textit{Mycobacterium avium} complex (MAC) \\
\end{longtable}

%============================================================
\section{診斷結果}
%============================================================

\subsection{HIV檢測}
\begin{itemize}
    \item HIV-1/HIV-2 antigen-antibody combination immunoassay：\textbf{Reactive}
    \item Confirmatory testing：\textbf{HIV-1 positive}
    \item \textbf{HIV-1 viral load：223,000 copies/mL}
    \item \textbf{CD4 count：33/$\mu$L} $\rightarrow$ Advanced HIV-1 infection
\end{itemize}

\subsection{\textit{Pneumocystis jirovecii} Pneumonia確診}
\begin{itemize}
    \item Induced sputum DFA for \textit{P. jirovecii}：Negative
    \item \textbf{Serum 1,3-$\beta$-D-glucan：>500 pg/mL}（參考值 <60）$\rightarrow$ 強烈陽性
    \item BAL fluid DFA for \textit{P. jirovecii}：\textbf{Positive}
\end{itemize}

\subsection{意外發現：\textit{Cryptococcus neoformans}}
\begin{itemize}
    \item 住院第3天BAL fluid fungal culture長出yeast colonies
    \item 鑑定為\textbf{\textit{Cryptococcus neoformans}}
    \item Serum cryptococcal antigen：Negative（可能因fungal burden低）
    \item Lumbar puncture排除CNS involvement
\end{itemize}

\keyfinding{
\textbf{最終診斷}：
\begin{enumerate}
    \item Advanced HIV-1 infection (CD4 33/$\mu$L)
    \item \textit{Pneumocystis jirovecii} pneumonia
    \item \textit{Cryptococcus neoformans} pneumonia
\end{enumerate}
}

%============================================================
\section{PJP診斷方法比較}
%============================================================

\begin{longtable}{p{4cm}p{5cm}p{5cm}}
\toprule
\textbf{檢測方法} & \textbf{敏感度} & \textbf{備註} \\
\midrule
\endfirsthead
\toprule
\textbf{檢測方法} & \textbf{敏感度} & \textbf{備註} \\
\midrule
\endhead
\bottomrule
\endfoot
Induced sputum DFA & 67--99\% (HIV+)\newline <50\% (immunocompetent) & 非侵入性，適合初步篩檢 \\
\textbf{BAL fluid DFA} & \textbf{90--99\% (HIV+)} & 敏感度最高，為gold standard \\
Serum 1,3-$\beta$-D-glucan & High sensitivity & 可早期提示，但非特異性 \\
Nucleic acid amplification & Highest sensitivity & 可能偵測到colonization，\newline 判讀需謹慎 \\
\end{longtable}

\note{\textit{P. jirovecii}現已歸類為ascomycetous fungus（真菌），而非原蟲。此菌無法在常規培養基上生長，因此immunofluorescence assay（如DFA）仍為標準診斷方法。}

%============================================================
\section{治療策略}
%============================================================

\subsection{PJP治療}

\subsubsection{Antimicrobial therapy}
\begin{itemize}
    \item \textbf{First-line：High-dose trimethoprim-sulfamethoxazole (TMP-SMX)}
    \begin{itemize}
        \item 療效優於其他替代藥物
        \item 需監測potassium及kidney function
        \item 療程3週
    \end{itemize}
    \item Alternative agents：Atovaquone, pentamidine, clindamycin + primaquine
\end{itemize}

\subsubsection{Adjunctive glucocorticoid therapy}
\textbf{適應症}（moderate-to-severe PJP）：
\begin{itemize}
    \item PaO$_2$ < 70 mmHg (room air)，或
    \item A-a gradient $\geq$ 35 mmHg
\end{itemize}

\begin{importantbox}
\textbf{本案病人SpO$_2$ 83\% (room air)}，屬clinically significant hypoxemia，符合glucocorticoid使用適應症。
\end{importantbox}

\subsection{Cryptococcal pneumonia治療}
\begin{itemize}
    \item Serum cryptococcal antigen陰性 $\rightarrow$ Localized disease
    \item Lumbar puncture排除CNS disease
    \item 治療：\textbf{Fluconazole}，總療程6個月
\end{itemize}

\subsection{Antiretroviral Therapy (ART)}

\subsubsection{ART啟動時機考量}
\begin{itemize}
    \item 早期啟動ART可減少AIDS-defining diseases及死亡風險
    \item 可增加viral suppression及patient retention in care
    \item \textbf{例外}：Cryptococcal meningitis患者應延遲ART（研究顯示早期啟動反而增加morbidity/mortality）
\end{itemize}

\subsubsection{本案處置}
\begin{itemize}
    \item CSF analysis排除cryptococcal meningitis
    \item 啟動\textbf{bictegravir-emtricitabine-tenofovir alafenamide}
\end{itemize}

%============================================================
\section{病程與預後}
%============================================================

\begin{itemize}
    \item 開始antimicrobial及glucocorticoid治療後症狀明顯改善
    \item 完成3週高劑量TMP-SMX後轉為prophylactic dose
    \item Fluconazole持續6個月
    \item \textbf{HIV viral load quickly became undetectable}
    \item 目前專注於HIV-informed health maintenance
    \item 病人獲知「\textbf{U = U: Undetectable = Untransmissible}」概念
\end{itemize}

%============================================================
\section{教學重點}
%============================================================

\subsection{老年人HIV篩檢的重要性}

\begin{itemize}
    \item 2021年美國HIV感染者中，\textbf{41\%為55歲以上}
    \item 55歲以上新診斷者中，\textbf{34\%診斷時已是AIDS}
    \item 老年人HIV篩檢率低於年輕族群
    \item 原因：
    \begin{itemize}
        \item 門診需討論的議題太多
        \item 社會觀感/尷尬
        \item 臨床醫師對老年人性行為的假設
    \end{itemize}
\end{itemize}

\subsection{CDC建議：性健康史「5 Ps」}
\begin{enumerate}
    \item \textbf{P}artners（性伴侶）
    \item \textbf{P}ractices（性行為方式）
    \item \textbf{P}rotection（STI預防措施）
    \item \textbf{P}ast（STI病史）
    \item \textbf{P}regnancy intention（懷孕意願）
\end{enumerate}

\subsection{本案關鍵教訓}

\begin{importantbox}
\begin{enumerate}
    \item \textbf{Herpes zoster可能暗示underlying immunodeficiency}
    \item \textbf{旅遊史及性行為史對任何年齡都重要}
    \item \textbf{Subacute respiratory syndrome + ground-glass opacities要考慮PJP}
    \item \textbf{1,3-$\beta$-D-glucan是PJP的有用biomarker}
    \item \textbf{Severely immunocompromised患者可能同時有多重伺機性感染}
    \item \textbf{早期ART啟動有益，但cryptococcal meningitis除外}
\end{enumerate}
\end{importantbox}

%============================================================
\section{考題}
%============================================================

\subsection{選擇題}

\begin{enumerate}
    \item 一位80歲女性，有8週的咳嗽、疲勞、體重減輕病史。5週前曾診斷herpes zoster。入院時SpO$_2$ 83\% (room air)，胸部CT顯示bilateral ground-glass opacities with central distribution and apical predominance。下列何者最可能是造成此影像表現的病因？
    \begin{enumerate}
        \item[(A)] Bacterial pneumonia
        \item[(B)] Pulmonary tuberculosis
        \item[(C)] \textbf{\textit{Pneumocystis jirovecii} pneumonia}
        \item[(D)] ANCA-associated vasculitis
        \item[(E)] Lymphangitic carcinomatosis
    \end{enumerate}
    \textit{答案：(C)。Central ground-glass opacities with apical predominance是PJP的典型CT表現。}

    \vspace{0.5cm}

    \item 關於\textit{Pneumocystis jirovecii} pneumonia的診斷，下列敘述何者\textbf{錯誤}？
    \begin{enumerate}
        \item[(A)] Induced sputum DFA在HIV感染者的敏感度為67--99\%
        \item[(B)] BAL fluid DFA是gold standard，敏感度可達90--99\%
        \item[(C)] Serum 1,3-$\beta$-D-glucan通常會升高
        \item[(D)] \textbf{\textit{P. jirovecii}可在常規培養基上生長}
        \item[(E)] Nucleic acid amplification test敏感度最高但可能偵測到colonization
    \end{enumerate}
    \textit{答案：(D)。\textit{P. jirovecii}是一種真菌，無法在常規培養基上生長，因此需要使用immunofluorescence assay進行診斷。}

    \vspace{0.5cm}

    \item 關於moderate-to-severe \textit{P. jirovecii} pneumonia的治療，下列敘述何者正確？
    \begin{enumerate}
        \item[(A)] Atovaquone是首選藥物
        \item[(B)] \textbf{High-dose trimethoprim-sulfamethoxazole是首選藥物}
        \item[(C)] 不需要使用adjunctive glucocorticoid
        \item[(D)] Glucocorticoid只用於PaO$_2$ < 50 mmHg的病人
        \item[(E)] 療程通常為7天
    \end{enumerate}
    \textit{答案：(B)。High-dose TMP-SMX是PJP的首選藥物。當PaO$_2$ < 70 mmHg或A-a gradient $\geq$ 35 mmHg時，應加上adjunctive glucocorticoid。療程為3週。}

    \vspace{0.5cm}

    \item HIV感染者的CD4 count降至何種程度時，較容易發生\textit{Pneumocystis jirovecii} pneumonia？
    \begin{enumerate}
        \item[(A)] $\leq$500/$\mu$L
        \item[(B)] $\leq$350/$\mu$L
        \item[(C)] \textbf{$\leq$200/$\mu$L}
        \item[(D)] $\leq$100/$\mu$L
        \item[(E)] $\leq$50/$\mu$L
    \end{enumerate}
    \textit{答案：(C)。PJP通常發生於CD4 count $\leq$200/$\mu$L的患者。這也是開始PJP prophylaxis的閾值。}

    \vspace{0.5cm}

    \item 關於HIV感染者開始antiretroviral therapy (ART)的時機，下列敘述何者正確？
    \begin{enumerate}
        \item[(A)] 所有伺機性感染都應延遲ART啟動
        \item[(B)] \textbf{Cryptococcal meningitis患者應延遲ART啟動}
        \item[(C)] PJP患者應延遲ART至完成3週治療後
        \item[(D)] CD4 count > 200/$\mu$L時才能開始ART
        \item[(E)] 需等待所有感染控制後才開始ART
    \end{enumerate}
    \textit{答案：(B)。早期啟動ART可減少AIDS progression及死亡，但研究顯示cryptococcal meningitis患者早期啟動ART反而增加morbidity/mortality，因此建議延遲ART。PJP患者則應early initiation of ART。}
\end{enumerate}

\subsection{問答題}

\begin{enumerate}
    \item 請說明為何本案80歲女性出現herpes zoster是一個重要的臨床線索？

    \textit{參考答案：Herpes zoster雖然常見於老年人（因年齡相關的免疫力下降），但也可能暗示underlying immunodeficiency。在本案中，herpes zoster的出現應促使臨床醫師考慮HIV感染的可能性，特別是當病人有相關風險因素時（如旅行史至HIV高盛行區、性行為史）。}

    \vspace{0.3cm}

    \item 為什麼本案的serum cryptococcal antigen是陰性，但BAL fluid culture卻長出\textit{Cryptococcus neoformans}？

    \textit{參考答案：Serum cryptococcal antigen testing對於disseminated或CNS cryptococcosis的敏感度較高，但在isolated pulmonary infection且fungal burden較低時可能呈陰性。因此，陰性的serum cryptococcal antigen不能排除localized pulmonary cryptococcosis。}

    \vspace{0.3cm}

    \item 請說明CDC建議的性健康史評估「5 Ps」包含哪些內容？為何這對老年人的HIV篩檢很重要？

    \textit{參考答案：5 Ps包括：Partners（性伴侶）、Practices（性行為方式）、Protection（STI預防措施）、Past（STI病史）、Pregnancy intention（懷孕意願）。這對老年人HIV篩檢很重要，因為2021年美國HIV感染者中41\%為55歲以上，且55歲以上新診斷者中34\%已是AIDS。老年人的HIV篩檢率偏低，部分原因是臨床醫師可能假設老年人沒有性行為。主動詢問性健康史可以識別需要HIV篩檢的高風險個體。}
\end{enumerate}

%============================================================
\section{參考文獻}
%============================================================

\begin{enumerate}
    \item Shrestha S, Brown T, Chu JT, Donnelly-Morell ML. Case 4-2026: An 80-Year-Old Woman with Cough and Hypoxemia. \href{https://doi.org/10.1056/NEJMcpc2513543}{\textit{N Engl J Med}. 2026;394:498-507.}

    \item Self WH, Balk RA, Grijalva CG, et al. Procalcitonin as a marker of etiology in adults hospitalized with community-acquired pneumonia. \href{https://doi.org/10.1093/cid/cix317}{\textit{Clin Infect Dis}. 2017;65:183-90.}

    \item Kanne JP, Yandow DR, Meyer CA. Pneumocystis jiroveci pneumonia: high-resolution CT findings in patients with and without HIV infection. \href{https://doi.org/10.2214/AJR.11.7329}{\textit{AJR Am J Roentgenol}. 2012;198:W555-W561.}

    \item Bozzette SA, Sattler FR, Chiu J, et al. A controlled trial of early adjunctive treatment with corticosteroids for Pneumocystis carinii pneumonia in the acquired immunodeficiency syndrome. \href{https://doi.org/10.1056/NEJM199011223232103}{\textit{N Engl J Med}. 1990;323:1451-7.}

    \item Boulware DR, Meya DB, Muzoora C, et al. Timing of antiretroviral therapy after diagnosis of cryptococcal meningitis. \href{https://doi.org/10.1056/NEJMoa1312884}{\textit{N Engl J Med}. 2014;370:2487-98.}

    \item Chang CC, Harrison TS, Bicanic TA, et al. Global guideline for the diagnosis and management of cryptococcosis. \href{https://doi.org/10.1016/S1473-3099(24)00127-3}{\textit{Lancet Infect Dis}. 2024;24(8):e495-e512.}
\end{enumerate}

\end{document}
